% CS 541 final report, on the course Overleaf template (same preamble, \lecture header, 6 sections).
% Overleaf layout: this file + preamble.tex (from the template) + images/ (upload report/figures/*.pdf)
%                  + tables/ (upload report/tables/*.tex).
% Numbers come from report/tables/*.tex and wiki/experiments/* (E1: wiki/experiments/e1_summary.md, because the
% E1 result files are on H's machine). Prose drafted by Claude from those sources at J's request; J and H must
% read, edit and be able to explain every sentence (disclosed in Section 6 and AI_USAGE.md).

\documentclass[11pt]{article}
\newtheorem{define}{Definition}

\usepackage{amsmath}
\usepackage{amsfonts}
\usepackage{pythonhighlight}
\usepackage{hyperref}
\usepackage{enumitem}
\usepackage{graphicx}
\graphicspath{ {./images/} }
\usepackage[linesnumbered,ruled,vlined, noline]{algorithm2e}

\SetKwInput{KwInput}{Input}
\SetKwInput{KwOutput}{Output}

\oddsidemargin=0.15in
\evensidemargin=0.15in
\topmargin=-.5in
\textheight=9in
\textwidth=6.25in
\usepackage[utf8]{inputenc}

\begin{document}
\input{preamble.tex}

% Arguments: {title}{date}{lecturer}{"By:" field}
\lecture{Design and Evaluation of a RAG Chatbot: Accuracy--Efficiency Trade-offs in Retrieval}{November 2026}
{Zhaozhuo Xu}
{Jing Cannon, Haiyan Zhu}

%==========================================================================================================
\section{What problem do we want to solve?}

Large language models (LLMs) write fluent answers, but in a specialised domain such as personal finance they
can state things that are not true, and the user has no way to check them. Retrieval-augmented generation
(RAG)~[2] addresses this: the system first retrieves documents that are relevant to the question and then asks
the LLM to answer \emph{only} from those documents and to cite them. The answer can then be no better than what
was retrieved. If the right document is missing from the retrieved set, the LLM either answers from the wrong
evidence or has to refuse.

This project studies the \textbf{retrieval layer} of a financial question-answering chatbot, not its user
interface. Retrieval is a nearest-neighbour search over text embeddings, and every design choice in it trades
accuracy against speed and memory. We ask three research questions:
\begin{itemize}[leftmargin=*]
\item \textbf{RQ1 --- search algorithm.} How much retrieval quality do we lose, and how much latency and
memory do we save, when we replace exact $k$-nearest-neighbour search with an approximate index (inverted file,
IVF)?
\item \textbf{RQ2 --- dimensionality.} How far can PCA reduce the 384-dimensional embeddings (to 128, 64, 32)
before quality drops too much?
\item \textbf{RQ3 --- reliability.} How do the number of retrieved documents $k$ and a similarity threshold
$\tau$ affect answer quality, and can $\tau$ make the system refuse questions that the corpus cannot answer?
\end{itemize}
The project uses several CS~541 topics directly: $k$-nearest neighbours with cosine similarity, PCA,
$k$-means clustering (the IVF index and our own clustering index), PyTorch, and an empirical analysis of
algorithmic efficiency.

%==========================================================================================================
\section{What datasets did you use?}

\paragraph{FiQA-2018.} We use the FiQA-2018 financial opinion QA dataset~[3] in the version distributed by the
BEIR benchmark~[5]. The corpus has \textbf{57,638} posts from finance forums (StackExchange
Investment and Reddit), and the questions are real user questions such as ``Is there ACH analogue in Asia?''.
Each question comes with relevance labels (qrels) that mark which posts answer it; all labels have score~1, and a
question has about 2.5 relevant posts on average (Table~\ref{tab:data}). One corpus is shared by all splits. We do
not train a model, so the train split is not used. All tuning is done on the 500 \textbf{dev} questions, and the
648 \textbf{test} questions were evaluated only once, after all settings were fixed (Section~\ref{sec:test}).
SciFact~[6] (5,183 abstracts, 300 questions) served only as a small smoke test of the pipeline.

The retrieval unit is a whole post, because the labels are given per post. Figure~\ref{fig:lengths} shows the
consequence: the median post has 120 word pieces, but our encoder reads at most 256, so \textbf{19\%} of the posts
are cut off before they are embedded. Splitting posts into chunks would avoid this, but it would no longer match
the post-level labels, so we kept whole posts and list the truncation as a limitation. Questions are short
(median 15 word pieces) and are never truncated.

\begin{table}[h]
\centering
\resizebox{\linewidth}{!}{\begin{tabular}{lllll}
\hline
Split & Queries & qrels pairs & Relevant docs / query (mean $\cdot$ median $\cdot$ max) & Used for \\
\hline
train & 5,500 & 14,166 & 2.58 $\cdot$ 2 $\cdot$ 23 & not used (no model is trained; PCA / IVF are fitted on the corpus) \\
dev & 500 & 1,238 & 2.48 $\cdot$ 2 $\cdot$ 20 & all tuning (E1--E3) \\
test & 648 & 1,706 & 2.63 $\cdot$ 2 $\cdot$ 15 & final evaluation, run once \\
SciFact test & 300 & 339 & --- & smoke test only (5,183 docs) \\
\hline
\end{tabular}
}
\caption{FiQA splits and their use. Relevant posts per question: mean, median, max.}
\label{tab:data}
\end{table}

\begin{figure}[h]
\centering
\includegraphics[width=0.9\linewidth]{data_lengths.pdf}
\caption{Length of FiQA texts in MiniLM word pieces. (a) The 57,638 corpus posts; the dotted line is the
encoder's 256-piece limit --- 19\% of posts are truncated before encoding. (b) The 500 dev questions.}
\label{fig:lengths}
\end{figure}

\paragraph{Unanswerable questions (for RQ3).} FiQA has no questions without an answer, but a refusal threshold can
only be evaluated with such questions. We therefore built them from dev: we picked 50 dev questions whose relevant
posts no other question needs and removed those posts (116 in total) from the corpus. Removing the labelled posts
does not guarantee that a question becomes unanswerable, because FiQA is a forum with many answers per topic and
the labels mark only some of them. We therefore checked every candidate by hand: for each one we read the top-5
retrieved posts and kept it only if none of them answered the question. ChatGPT drafted a verdict and a reason for
each candidate; we discussed every candidate together and made the final decision. Only \textbf{10 of the 50}
candidates were really unanswerable; in the other 40, an unlabelled post still answered the question. These 10
questions, together with the 450 remaining answerable dev questions, are the refusal benchmark.

%==========================================================================================================
\section{What models have you tried?}

\paragraph{Pipeline.} Our system has four steps: (1) embed the question, (2) retrieve the $k$ most similar posts,
(3) refuse if even the best post is not similar enough, (4) otherwise let an LLM answer from the $k$ posts with
citations (Algorithm~\ref{alg:rag}).

\paragraph{Encoder.} All texts are embedded with \texttt{sentence-transformers/all-MiniLM-L6-v2}~[4], a small
Sentence-BERT model that maps a text to a 384-dimensional vector. We L2-normalise every vector, so cosine
similarity equals the inner product. The encoder is fixed in every experiment; only the retrieval side changes.

\paragraph{Search back ends (RQ1).} All back ends search the same normalised embeddings:
\begin{enumerate}[label=(\alph*), leftmargin=*]
\item \textbf{TorchKNN} --- our exact brute-force search in PyTorch: one matrix product between the query and
all 57,638 document vectors, then \texttt{topk}. Cost $O(N d)$ per query.
\item \textbf{FAISS Flat}~[1] (\texttt{IndexFlatIP}) --- the same exact search in an optimised C++ library, used
as a correctness reference for (a).
\item \textbf{FAISS IVF} (\texttt{IndexIVFFlat}) --- approximate search. $k$-means splits the corpus into
$n_{\text{list}} = 240 \approx \sqrt{57{,}638}$ cells; a query is compared only with the documents in its
$n_{\text{probe}}$ closest cells. We tried $n_{\text{probe}} \in \{1, 4, 16, 64\}$.
\item \textbf{Our $k$-means index} (bonus) --- the IVF idea implemented by us in PyTorch (20 Lloyd iterations,
240 clusters, inverted lists), to compare our own implementation with FAISS.
\end{enumerate}

\paragraph{PCA (RQ2).} PCA is fitted on the corpus embeddings only and reduces them to 128, 64 or 32 dimensions.
The reduced vectors are normalised again, and every query goes through exactly the same transform.

\paragraph{Generator and refusal (RQ3).} The LLM is \texttt{openai/gpt-oss-120b}, an open-weight model, called
through the Groq API with temperature~0, low reasoning effort and at most 400 output tokens. The prompt gives the
$k$ posts with their ids and tells the model to answer only from them, to cite the id after every claim, and
otherwise to reply with one fixed refusal sentence. The proposal planned a small local instruction-tuned model; we
switched to a larger hosted open-weight model because we expected small models to be weak at citing and staying
inside the documents, and because a free API with published limits needed no GPU and no multi-gigabyte download. Before
the LLM is called, the system refuses on its own if the top-1 cosine similarity is below a threshold $\tau$.

\begin{algorithm}[h]
\caption{Answer or refuse}
\label{alg:rag}
\KwInput{question $q$, document embeddings $D$, index, $k$, threshold $\tau$}
\KwOutput{an answer with cited post ids, or a refusal}
$e \leftarrow \text{normalise}(\text{PCA}(\text{MiniLM}(q)))$ \tcp*{PCA only in config B}
$(s_1, \dots, s_k), (d_1, \dots, d_k) \leftarrow \text{index.search}(e, k)$ \tcp*{cosine scores, best first}
\eIf{$s_1 < \tau$}{\Return refusal \tcp*{no post is similar enough}}
{\Return LLM(prompt$(q, d_1, \dots, d_k)$) \tcp*{may still return the refusal sentence}}
\end{algorithm}

\paragraph{Final configurations.} From the dev results we fixed two configurations before touching the test set:
\textbf{A (accurate)} = 384-d exact TorchKNN, $\tau = 0.60$; and \textbf{B (efficient)} = PCA 128-d + FAISS IVF
($n_{\text{list}}=240$, $n_{\text{probe}}=16$), $\tau = 0.66$. Generation uses A with $k=5$. A Streamlit demo runs
the same pipeline with config A.

%==========================================================================================================
\section{How to evaluate the performance of the model on your dataset?}

\paragraph{Retrieval quality.} For each question we compare the top-10 posts with the labelled relevant posts:
\begin{itemize}[leftmargin=*]
\item \textbf{Recall@$k$} ($k = 5, 10$): the fraction of the relevant posts that appear in the top $k$. Recall@10 is
our main metric. Recall@5 matters because only the top 5 posts are given to the LLM.
\item \textbf{MRR@10}: the mean of $1/\text{rank}$ of the first relevant post (0 if none is in the top 10).
\item \textbf{nDCG@10}: like recall, but a relevant post counts more the higher it is ranked
($\text{DCG} = \sum_i \text{rel}_i / \log_2(i+1)$, divided by the best possible DCG).
\end{itemize}
Rank matters for two reasons: the LLM only sees the first five posts, and the refusal rule only looks at the
top-1 score.

\paragraph{Efficiency.} Latency is measured per single query (batch size 1, as in a chatbot), after a warm-up, over
three passes of all queries; we report the mean and the 95th percentile (p95). We also report the index memory.
Latencies are only compared within one session on one machine; differences below about 10\% are within the noise
we observed and are treated as ties.

\paragraph{Refusal.} On the 450 answerable and 10 unanswerable dev questions we measure the \emph{answer rate}
(answerable questions that are answered) and the \emph{correct-refusal rate} (unanswerable questions that are
refused). We choose $\tau$ by \emph{balanced accuracy}, the mean of the two rates. Plain accuracy would be
misleading: with 450 answerable and 10 unanswerable questions, ``always answer'' already reaches 97.8\% accuracy
but only 0.5 balanced accuracy.

\paragraph{Answer quality.} The LLM's answers are scored on a 1--3 rubric with two separate scores.
\emph{Correctness} compares the answer with the labelled relevant posts (3 = covers the main point without errors,
2 = partly right, 1 = wrong). \emph{Groundedness} checks every claim against the $k$ posts the model was given
(3 = everything is supported and correctly cited, 2 = mostly supported, 1 = made up or from outside knowledge).
The two scores fail in different ways: ``correct but not grounded'' means the model used its own knowledge, which
the user cannot verify; ``grounded but wrong'' means retrieval gave the wrong posts. Refusals are not scored; they
are labelled \emph{refused} (the posts really lack the answer) or \emph{refused wrongly} (they contain it). ChatGPT
drafted a score and a reason for every answer; we discussed each item together and decided the final scores. This
is one joint score per answer, so we cannot report agreement between two independent raters.

\paragraph{Protocol.} Every setting ($n_{\text{probe}}$, PCA dimension, $k$, $\tau$) was chosen on dev. The two
final configurations were then written down and evaluated \textbf{once} on the 648 test questions; nothing was
changed afterwards. For generation we sampled 100 test questions with a fixed random seed (42), which is also used
everywhere else.

%==========================================================================================================
\section{How does your model perform?}

\subsection{E1: exact vs.\ approximate search (RQ1)}

Table~\ref{tab:e1} summarises E1 on dev. Our TorchKNN and FAISS Flat give identical scores on all four quality
metrics, which validates our implementation against a reference. Their latencies (0.77 vs.\ 0.83 ms p95) are a tie:
both do the same $O(Nd)$ scan, which for a single query is limited by memory bandwidth.

IVF shows a clear accuracy--speed curve. With $n_{\text{probe}} = 1$ it is 47$\times$ faster but loses
18.35 points of Recall@10; each larger $n_{\text{probe}}$ buys accuracy back at a higher price in speed. At
$n_{\text{probe}} = 64$ (about 27\% of the corpus visited) IVF reproduces the exact result, so the losses at small
$n_{\text{probe}}$ come from relevant posts that lie in cells the query does not visit. We chose
$n_{\text{probe}} = 16$ as the trade-off: $-2.82$ points for a 6.7$\times$ lower p95. The index adds less than
1~MB (centroids and inverted lists), so IVF saves time but not memory.

Our own $k$-means index reaches the same accuracy as FAISS IVF at the same settings (Recall@10 0.4435 vs.\ 0.4383)
but is 17$\times$ slower than FAISS IVF and even slower than exact search, because it loops over queries and
cells in Python while FAISS uses optimised C++. Our implementation reproduces the accuracy of an industrial index,
but not its speed.

To check that the speed-up of IVF grows with the corpus, we repeated exact search and IVF ($n_{\text{list}} = 240$,
$n_{\text{probe}} = 16$) on random sub-corpora of 10,000 and 25,000 posts that keep every relevant post
(Figure~\ref{fig:scaling}, one thread for both methods). The speed-up grows from 3.4$\times$ to 6.5$\times$ and
7.2$\times$ on the full corpus, while the loss stays at 2.5--2.8 points of Recall@10. It stays below the bound of
$240/16 = 15\times$ because IVF also scans a fixed share of the corpus and pays fixed costs per query. Recall@10
falls with $N$ for both methods only because a larger corpus adds distractors.

\begin{figure}[h]
\centering
\includegraphics[width=0.9\linewidth]{e1_scaling.pdf}
\caption{Exact search vs.\ FAISS IVF ($n_{\text{probe}} = 16$) on FiQA dev for three corpus sizes, one thread.
Left: p95 latency per query (log scale); right: speed-up of IVF over exact search, with the bound
$240/16 = 15\times$ dashed.}
\label{fig:scaling}
\end{figure}

\begin{table}[h]
\centering
\resizebox{\linewidth}{!}{\begin{tabular}{llllll}
\hline
Method & Recall@10 & $\Delta$ vs exact (pt) & nDCG@10 & p95 (ms) & Speed-up \\
\hline
TorchKNN (ours, exact) & 0.4665 & +0.00 & 0.3881 & 0.765 & 1.0$\times$ \\
FAISS Flat (exact) & 0.4665 & +0.00 & 0.3881 & 0.829 & 0.9$\times$ \\
FAISS IVF, nprobe 1 & 0.2831 & -18.35 & 0.2387 & 0.016 & 46.6$\times$ \\
FAISS IVF, nprobe 4 & 0.3734 & -9.32 & 0.3063 & 0.038 & 20.2$\times$ \\
FAISS IVF, nprobe 16 & 0.4383 & -2.82 & 0.3631 & 0.114 & 6.7$\times$ \\
FAISS IVF, nprobe 64 & 0.4665 & +0.00 & 0.3882 & 0.391 & 2.0$\times$ \\
K-means index (ours), nprobe 16 & 0.4435 & -2.30 & 0.3674 & 1.978$^{*}$ & 0.4$\times$ \\
\hline
\end{tabular}
}
\caption{E1 on FiQA dev (500 questions, 384-d, full corpus). IVF uses $n_{\text{list}} = 240$. p95 latency per
single query at the library's natural thread count; speed-up relative to TorchKNN. $^{*}$Measured with one
thread.}
\label{tab:e1}
\end{table}

\subsection{E2: PCA dimensionality (RQ2)}

Table~\ref{tab:pca} and Figure~\ref{fig:pca} show the effect of PCA with exact search, so any change is caused by
PCA alone. Memory falls in proportion to the dimension. Quality falls more slowly than the explained variance at
first, but the gap closes quickly: 128 dimensions keep 79.5\% of the variance and lose only 2.54 points of
Recall@10, while 64 dimensions lose 9.94 points and 32 lose 19.79. The first discarded directions are mostly noise
for retrieval; below 128 PCA starts to remove directions that carry the ranking signal. The speed-up is smaller
than the reduction in dimension (3$\times$ fewer dimensions give 2.6$\times$ lower p95) because part of the cost
per query does not depend on $d$. We therefore chose \textbf{128 dimensions}: 3$\times$ less memory and
2.6$\times$ faster for 2.54 points.

An important difference from IVF: the IVF loss can be bought back with a larger $n_{\text{probe}}$, while the
information removed by PCA is gone. Combining both in config B (128-d + IVF 16) costs 4.88 points of Recall@10 on
dev, slightly less than the sum of the two separate losses, and is 30$\times$ faster than config A in the same
session (Table~\ref{tab:final}).

\begin{figure}[h]
\centering
\includegraphics[width=0.45\linewidth]{e2_variance.pdf}
\includegraphics[width=0.9\linewidth]{e2_pca_tradeoff.pdf}
\caption{E2 on FiQA dev. Top: cumulative explained variance of PCA. Bottom: (a) Recall@10 and nDCG@10 vs.
dimension; (b) Recall@10 vs.\ p95 latency.}
\label{fig:pca}
\end{figure}

\begin{table}[h]
\centering
\resizebox{\linewidth}{!}{\begin{tabular}{llllllll}
\hline
Dim & Variance kept & Recall@10 & $\Delta$ vs 384-d (pt) & nDCG@10 & p95 (ms) & Speed-up & Memory (MB) \\
\hline
384 & 1.000 & 0.4665 & +0.00 & 0.3881 & 2.567 & 1.0$\times$ & 84.43 \\
128 & 0.795 & 0.4412 & -2.54 & 0.3586 & 0.975 & 2.6$\times$ & 28.14 \\
64 & 0.614 & 0.3671 & -9.94 & 0.2937 & 0.663 & 3.9$\times$ & 14.07 \\
32 & 0.456 & 0.2686 & -19.79 & 0.2079 & 0.464 & 5.5$\times$ & 7.04 \\
\hline
\end{tabular}
}
\caption{E2: PCA dimension vs.\ quality, latency and memory (FiQA dev, exact search). Latencies are from a
different machine than E1 and must not be compared with Table~\ref{tab:e1}.}
\label{tab:pca}
\end{table}

\subsection{E3: refusal threshold and top-$k$ (RQ3)}

\paragraph{Threshold $\tau$.} The top-1 cosine similarity separates answerable from unanswerable questions
reasonably well (Figure~\ref{fig:tau}, top): the medians are 0.652 and 0.549, and the area under the ROC curve is
0.823. Before our manual check, with all 50 candidates, the AUC was only 0.56, because most ``unanswerable''
candidates were in fact answerable. Without the check, threshold refusal would have looked useless. The grid search
(Figure~\ref{fig:tau}, bottom; Table~\ref{tab:tau}) gives $\tau = 0.60$ for config A: it answers 76\% of the
answerable questions and refuses 8 of the 10 unanswerable ones (balanced accuracy 0.78). Config B needs its own
threshold, $\tau = 0.66$, because PCA raises all cosine values. With only 10 unanswerable questions this result is
uncertain: the 95\% confidence interval for 8/10 is roughly 0.49--0.94.

\begin{table}[h]
\centering
\resizebox{\linewidth}{!}{\begin{tabular}{lllllll}
\hline
Config & Median top-1 (answerable / unanswerable) & AUC & $\tau$ & Answer rate & Correct refusal & Balanced \\
\hline
A --- 384-d exact & 0.652 / 0.549 & 0.823 & 0.60 & 0.760 & 0.80 & 0.780 \\
B --- 128-d PCA + IVF 16 & 0.680 / 0.574 & 0.826 & 0.66 & 0.609 & 0.90 & 0.754 \\
\hline
\end{tabular}
}
\caption{E3: refusal threshold on FiQA dev (450 answerable, 10 verified unanswerable questions).}
\label{tab:tau}
\end{table}

\begin{figure}[h]
\centering
\includegraphics[width=0.6\linewidth]{e3_top1_distribution.pdf}
\includegraphics[width=0.9\linewidth]{e3_tau_curve.pdf}
\caption{E3 on FiQA dev. Top: top-1 cosine of answerable vs.\ verified unanswerable questions. Bottom: refusal
threshold grid search for configs A and B.}
\label{fig:tau}
\end{figure}

\paragraph{Number of documents $k$.} On 20 dev questions (config A), more context helps correctness and makes the
LLM refuse less often, but it lowers groundedness (Figure~\ref{fig:k}): from $k=1$ to $k=10$ the number of answered
questions rises from 11 to 17 and correctness from 2.36 to 2.65, while groundedness falls from 3.00 to 2.65. With
more posts the model combines more sources and makes more claims that are not exactly supported by them. We chose
$k=5$ as the compromise (correctness 2.53, groundedness 2.87). With 20 questions these differences are small and
should be read as trends.

\begin{figure}[h]
\centering
\includegraphics[width=0.9\linewidth]{e3_quality_vs_k.pdf}
\caption{Answer quality vs.\ $k$ on 20 FiQA dev questions (config A).}
\label{fig:k}
\end{figure}

\subsection{Final test evaluation (run once)}
\label{sec:test}

\paragraph{Retrieval.} The dev conclusions hold on the test set (Table~\ref{tab:final}, Figure~\ref{fig:test},
top). Both configurations score slightly lower on test than on dev (Recall@10: A 0.4665 $\rightarrow$ 0.4413,
B 0.4178 $\rightarrow$ 0.3878), as expected after choosing settings on dev. The gap between them stays almost the
same ($-5.35$ points on test, $-4.88$ on dev), and B needs 3$\times$ less memory. During the test run the machine
was under heavy background load, and A's multithreaded search slowed down (p95 11.6 ms vs.\ 2.7 ms on dev) while
B's small search did not. The test latency of A is therefore not reliable, and we report the dev same-session
speed-up of \textbf{30$\times$} as the speed result. We did not re-run the test, by protocol.

\begin{table}[h]
\centering
\resizebox{\linewidth}{!}{\begin{tabular}{lllllllll}
\hline
Config & Split & Queries & Recall@5 & Recall@10 & nDCG@10 & MRR@10 & p95 (ms) & Memory (MB) \\
\hline
A --- 384-d exact & dev & 500 & 0.3797 & 0.4665 & 0.3881 & 0.4532 & 2.700 & 84.43 \\
A --- 384-d exact & test & 648 & 0.3671 & 0.4413 & 0.3687 & 0.4451 & 11.630 & 84.43 \\
B --- 128-d PCA + IVF 16 & dev & 500 & 0.3343 & 0.4178 & 0.3388 & 0.3956 & 0.089 & 28.14 \\
B --- 128-d PCA + IVF 16 & test & 648 & 0.3258 & 0.3878 & 0.3275 & 0.4060 & 0.080 & 28.14 \\
\hline
\end{tabular}
}
\caption{Final configurations on FiQA dev and test. Test latency of A was inflated by background load.}
\label{tab:final}
\end{table}

\paragraph{Generation.} Of the 100 test questions, the threshold refused 16, so the system's answer rate is 0.84
(0.76 on dev). Of the remaining 84, the LLM itself refused 14: in 10 cases correctly, in 4 cases although the
answer was in the posts. The system gave a real answer to \textbf{70} questions, with mean correctness 2.26 and
groundedness 2.36 (Table~\ref{tab:gen}). These scores are lower than on dev (2.53 / 2.87 on 20 questions); the
dev sample is small, so part of this difference may be chance.

The most useful result is the split by retrieval success (Figure~\ref{fig:test}, bottom b). When a labelled
relevant post is among the five posts given to the LLM, correctness is 2.36; when none is, it drops to 1.87, while
groundedness stays the same (2.36 vs.\ 2.33). The LLM stays faithful to the posts it receives even when they are
the wrong posts. In our system, errors come mainly from \textbf{retrieval}, not from generation, so better retrieval
would help more than a better prompt. We also saw that 13 of the answered questions hit the 400-token output limit
(the reasoning tokens of gpt-oss count toward it), and these answers scored lower (2.08 / 2.15).

\begin{table}[h]
\centering
\resizebox{0.75\linewidth}{!}{\begin{tabular}{llll}
\hline
Subset & n & Correctness & Groundedness \\
\hline
Answered by the system & 70 & 2.26 & 2.36 \\
\ldots{} relevant doc in top 5 & 55 & 2.36 & 2.36 \\
\ldots{} no relevant doc in top 5 & 15 & 1.87 & 2.33 \\
\ldots{} answer hit the 400-token cap & 13 & 2.08 & 2.15 \\
$\tau$-refused, LLM answered anyway & 7 & 2.29 & 2.57 \\
Dev reference (k = 5, 20 questions) & 20 & 2.53 & 2.87 \\
\hline
\end{tabular}
}
\caption{Answer quality on 100 FiQA test questions (config A, $k=5$, $\tau=0.60$), mean scores 1--3. The
``$\tau$-refused'' row scores answers the LLM gave for questions the system refused; they are not shown to the
user.}
\label{tab:gen}
\end{table}

\begin{figure}[h]
\centering
\includegraphics[width=0.9\linewidth]{final_dev_vs_test.pdf}
\includegraphics[width=0.9\linewidth]{final_generation.pdf}
\caption{Top: retrieval quality of A and B on dev vs.\ test. Bottom: the final system on 100 test questions ---
(a) outcome per question, (b) answer quality with and without a relevant post in the top 5.}
\label{fig:test}
\end{figure}

\subsection{Answers to the research questions}

\paragraph{RQ1 --- How much retrieval quality do we trade for faster search?}
On FiQA dev, FAISS IVF with 16 probes reduced p95 retrieval latency from 0.765 ms to 0.114 ms, a 6.7$\times$
speed-up. Recall@10 fell from 46.65\% to 43.83\%, a drop of 2.82 percentage points, and index memory stayed about
the same. We chose 16 probes as the operating point because it retained most of the exact-search recall while
substantially reducing query latency.

\paragraph{RQ2 --- How far can we reduce the embedding dimension?}
Among the dimensions we tested, 128 offered the best trade-off. Reducing the embeddings from 384 to 128 dimensions
cut embedding memory to one third and made exact search 2.6$\times$ faster at p95, with a 2.54-point drop in
Recall@10. The recall losses grew to 9.94 points at 64 dimensions and 19.79 points at 32. We therefore selected
128 dimensions for the smaller configuration.

\paragraph{RQ3 --- How do the number of retrieved posts and the similarity threshold affect reliability?}
On 20 dev questions, increasing $k$ from 1 to 10 raised mean correctness from 2.36 to 2.65 and reduced mean
groundedness from 3.00 to 2.65 among the answers the model gave. The number of answers also increased from 11 to
17, so the averages reflect different sets of questions. We chose $k=5$ as a compromise. With $\tau=0.60$,
configuration A passed 76\% of answerable dev questions through the similarity gate and rejected 8 of the 10
manually reviewed unanswerable cases. The threshold was selected on this same small dev benchmark, so its
performance on unseen unanswerable questions remains uncertain.

\subsection{Limitations}
\begin{itemize}[leftmargin=*]
\item 19\% of the posts are truncated by the encoder, so the end of a long answer is invisible to retrieval.
\item The threshold was chosen and evaluated on the same 10 unanswerable dev questions, and there is no
unanswerable test set. The correct-refusal rate is therefore uncertain and probably optimistic.
\item Answer scores are one joint, ChatGPT-assisted judgement per answer. The proposal planned 30 answers scored
independently by both of us; we did not do this, so we have no inter-rater agreement.
\item The choice of $k$ rests on only 20 dev questions.
\item 13 of the 100 test answers hit the 400-token limit.
\item Test latency of config A is unreliable because of background load; E1 and E2 latencies come from different
machines and are only compared within one table.
\item The LLM is a hosted model that can change over time; we recorded the model id and date with every answer.
\end{itemize}

%==========================================================================================================
\section{How does your AI teammate perform?}
% Written by J and H; Claude added disclosure details (code scope, report/slides drafting, ChatGPT verdicts, final decision by J+H).

Claude Code wrote most of the code at our request and helped test it, ran experiments, and turned the results into
tables and figures. The exception is the search back ends (exact KNN, FAISS wrappers and our $k$-means index):
Haiyan implemented them, and Claude only explained concepts and helped locate bugs. Claude also drafted the text of
this report and the presentation, which we then revised. ChatGPT drafted the answer scores (dev and test) and the
verdicts for the 50 unanswerable candidates, and gave a written reason for each judgment; we discussed every item
and decided the final judgments together. The written reasons made it easier to review why a response received a
particular score. Every use of AI is listed in \texttt{AI\_USAGE.md} in our repository.

The AI sometimes made confident claims that we had to check. For example, Claude said a merge had been pushed, but
the repository status showed otherwise. In the scoring work, an early ChatGPT draft assigned full groundedness credit to
refusals. We corrected that by leaving both quality scores blank for refusals and counting them separately. We
checked the code with tests and used the saved run records to verify the results. The scoring process produced one
set of judgments, so we cannot estimate agreement between independent raters.

%==========================================================================================================
\clearpage
\section*{References}
\begin{enumerate}[leftmargin=*]
\item Johnson, J., Douze, M., \& J\'egou, H. (2019). Billion-scale similarity search with GPUs. \emph{IEEE Transactions on Big Data}.
\item Lewis, P., et al. (2020). Retrieval-Augmented Generation for Knowledge-Intensive NLP Tasks. \emph{NeurIPS}.
\item Maia, M., et al. (2018). WWW'18 Open Challenge: Financial Opinion Mining and Question Answering. \emph{Companion Proc.\ of The Web Conference}.
\item Reimers, N., \& Gurevych, I. (2019). Sentence-BERT: Sentence Embeddings using Siamese BERT-Networks. \emph{EMNLP}.
\item Thakur, N., et al. (2021). BEIR: A Heterogeneous Benchmark for Zero-shot Evaluation of Information Retrieval Models. \emph{NeurIPS Datasets and Benchmarks}.
\item Wadden, D., et al. (2020). Fact or Fiction: Verifying Scientific Claims. \emph{EMNLP}.
\end{enumerate}

\end{document}
